\section{Notation}\label{apdx:notation}

We use $\|x\|$ to denote the Euclidean norm of a vector $x$, $\|A\|_\text{op}$ for the spectral norm and $\|A\|_F$ for the Frobenius norm of a matrix $A$.
$x_i$ denotes the $i$th element of a vector $x$, while $x^{(i)}$ denotes the $i$th vector in a list of vectors.
We denote $\text{nnz}(\cdot)$ the number of non-zero entries in a vector or matrix. For two functions $f, g$ we say $f(z) = O(g(z))$ if there is a constant $C$, such that $0 \le f(z) \le C g(z)$ for $z$ large enough. Similarly, $f(z) = \Omega(g(z))$ if there is a constant $C$, such that $f(z) \ge C g(z)$ for $z$ large enough. $\tilde O$ and $\tilde \Omega$ omit poly-logarithmic factors. 

Throughout the paper, we use $d$ for the dimension of input vectors, $n$ for the number of samples in the dataset, and $M$ for the dimension of the sketch or random features.  